\documentclass[12pt]{article}
\usepackage[utf8]{inputenc}
\usepackage[T1]{fontenc}
\usepackage{hyperref}
\title{Seta Language Resources for K-12 Curriculum Development}
\author{Compiled by Claude Sonnet 5.5 (claude-sonnet-5-5) for Andr\'as Kornai}
\date{}
\begin{document}
\maketitle

\section{Introduction}
Seta (ISO 639-3: stf; Glottocode seta1245) is a Torricelli language of the One (West Wapei) branch, spoken in Sandaun Province, Papua New Guinea. Wikipedia cites 180 speakers from the 2000 census, while Joshua Project gives about 300 based on Ethnologue 2016. Searches found no dictionary, grammar, primer or school materials. Seta should not be confused with Seti, which has a separate Wikipedia entry.

\section{Primary Educational Resources}
None were found.

\section{Linguistic and Reference Materials}
\begin{itemize}
\item Wikipedia, ``Seta language'': \url{https://en.wikipedia.org/wiki/Seta_language}
\item Wikipedia, ``One languages'' (the branch Seta belongs to): \url{https://en.wikipedia.org/wiki/One_languages}
\item Wikipedia, ``Torricelli languages'': \url{https://en.wikipedia.org/wiki/Torricelli_languages}
\item ASJP word list, Seta (compiled by Matthew S. Dryer; it cites a 2001 University of Sydney MA thesis by Crowther): \url{https://asjp.clld.org/languages/SETA_PLEMA}
\item Endangered Languages Project, bibliography for Seta: \url{https://endangeredlanguages.com/lang/5421/bibliography?hl=pt-br}
\end{itemize}

\section{Religious and Community Resources}
\begin{itemize}
\item Joshua Project, Seta in Papua New Guinea (it lists Bible translation status as ``Translation Needed''): \url{https://joshuaproject.net/people_groups/14874}
\end{itemize}

\section{Digital and Audio Resources}
No audio, video or corpus resources were found beyond the ASJP word list above.

\section{Curriculum Development Context}
The language is endangered, and not all young people speak it. No known school use exists. Any curriculum would have to begin with documentation led by the community, and it would draw on the lingua francas of the area. The small word list and the thesis cited by ASJP are the only linguistic starting points identified here.
\end{document}